\selectlanguage{ngerman}%
\frenchspacing%
\begin{ledgroupsized}[r]{120mm}%
\footnotesize%
\pstart%
\noindent\textbf{Überlieferung:}%
\pend%
\end{ledgroupsized}%
\begin{ledgroupsized}[r]{114mm}%
\footnotesize%
\pstart%
\parindent -6mm%
\makebox[6mm][l]{\textit{L}}%
Aufzeichnung:
LH XXXV\;11,\;14 Bl.\;15.
Ein Blatt 2\textsuperscript{o};
Ränder beschnitten. %
Zwei ganzseitig beschriebene Seiten;
Textfolge:
Bl.\;\href{\Scan{LH_35_11_14_015v}}{15\,v\textsuperscript{o}}, \href{\Scan{LH_35_11_14_015r}}{15\,r\textsuperscript{o}}.
Teilentwürfe zum Diagramm \bracks{\textit{Fig.\:3}} auf Bl.\;\href{\Scan{LH_35_11_14_015v}}{15\,v\textsuperscript{o}}.%
\pend%
\end{ledgroupsized}%
\vspace{5mm}%
\begin{ledgroup}%
\footnotesize%
\pstart%
\noindent%
\textbf{Datierungsgründe:}
Die Datierung der Aufzeichnung N.\;\ref{58943} beruht auf %
inhaltlichen Betrachtungen.
Siehe die ausführliche Begründung auf S.\;\pageref{58943_dat}.%
\pend%
\end{ledgroup}%
\selectlanguage{latin}%
\frenchspacing%
\vspace{8mm}%
\normalsize%
\pstart%
\noindent%
\bracks{%
\href{\Scan{LH_35_11_14_015v}}%
{15\,v\textsuperscript{o}}%
}%
\hfill%
\edtext{%
De Tensionibus funium%
\indexS{tensio funis}%
\indexS{funis tensus}
\edtext{%
et obliquis directionibus%
\indexS{directio obliqua}%
}{%
\lemma{et}%
\Bfootnote{%
\textit{(1)}~potentiis\indexS{potentia} ob %
\textit{(2)}~obliquis directionibus%
~\textit{L}%
}}%
}{%
\lemma{\textit{Am Rand:}}%
\Afootnote{%
Hic specimen Methodi%
\indexS{specimen methodi de maximo}
de Maximo%
\indexS{methodus de maximo eligendo}%
\indexS{maximum eligendum ex infinitis maximis}
ex infinitis maximis eligendo;
seu cum infinities infinitae varietates%
\indexS{varietas infinities infinita}
per solam maximae conditionem determinantur.%
}}%
\hfill%
\phantom{\bracks{15\,v\textsuperscript{o}}}%
\pend%
\vspace{0.5em}%
\pstart%
\noindent%
Quanta
\edtext{%
vi \textit{A} et \textit{B} sustentacula%
\indexS{sustentaculum}
trahantur a pondere \textit{C};%
}{%
\lemma{vi}%
\Bfootnote{%
\textit{(1)}~pondus \textit{A} %
\textit{(a)}~\angles{--} %
\textit{(b)}~et %
\textit{(2)}~\textit{A} et \textit{B} %
\bracks{...} pondere\:\textit{C};%
~\textit{L}%
}%
\lemma{\textit{A} et \textit{B} sustentacula}%
\Cfootnote{%
Vgl. das Diagramm \bracks{\textit{Fig.\:1}}.%
}%
}
seu suspensis ponderibus \textit{D} et \textit{E}%
\indexS{pondus suspensum}
quaeritur ratio eorum,%
\indexS{ratio ponderum suspensorum}
ut omnia sint in aequilibrio.%
\indexS{pondera in aequilibrio}%
\indexS{aequilibrium ponderum suspensorum}
\pend%
\pstart%
Sit
\edtext{%
\textit{D} et \textit{E} ponderum%
}{%
\lemma{\textit{D} et \textit{E} ponderum}%
\Cfootnote{%
Vgl. das Diagramm \bracks{\textit{Fig.\:2}}.%
}}
centrum grav. \textit{F} in recta \textit{DE},
et centrum grav. omnium \textit{G} in recta \textit{CF}.%
\indexS{centrum gravitatis ponderum suspensorum}
Si jam ponamus
\textit{C} ascendere in \textit{(C)}%
\indexS{pondus ascendens}%
\indexS{ascensus ponderis}
descendentibus \textit{D} et \textit{E}
in \textit{(D)} et \textit{(E)},%
\indexS{pondus descendens}%
\indexS{descensus ponderis}
et omnia
\edtext{%
prius in aequilibrio
fu\bracks{iss}e,
tunc si manebunt%
}{%
\lemma{prius}%
\Bfootnote{%
in aequilibrio %
\bars{fuere \textit{ändert Hrsg.}}, %
\bars{tunc si \textit{erg.}} manebunt%
~\textit{L}%
}}
nunc quoque in aequilibrio,%
\indexS{pondera in aequilibrio}%
\indexS{aequilibrium ponderum suspensorum}
et punctum \textit{G} manebit immotum%
\bracks{,}
puncta autem \textit{(C)}, \textit{G}, \textit{(F)} erunt in eadem recta.
Et debent esse
elevatio ipsius \textit{C}%
\indexS{elevatio centri gravitatis}
et depressio ipsius \textit{F}%
\indexS{depressio centri gravitatis}
reciproce ut \textit{C} ad
\edtext{%
$ \overline{D + E} $.
Verum haec assertio non est certa.%
\indexS{assertio certa}
Videtur enim%
}{%
\lemma{$ \overline{D + E} $.}%
\Bfootnote{%
\textit{(1)}~Ita enim %
\textit{(2)}~Verum haec %
\bracks{...} Videtur enim %
~\textit{L}%
}}
potius centrum gravitatis%
\indexS{centrum gravitatis ponderum suspensorum}
attolli,
sive \textit{C} ascendat%
\indexS{pondus ascendens}%
\indexS{ascensus ponderis}
sive descendat.%
\indexS{pondus descendens}%
\indexS{descensus ponderis}%
\pend%
\pstart%
Sint tria
\edtext{%
pondera \textit{C}, \textit{D}, \textit{E},%
}{%
\lemma{pondera \textit{C}, \textit{D}, \textit{E}}%
\Cfootnote{%
Vgl. das Diagramm \bracks{\textit{Fig.\:3}}.%
}}
quorum
duo \textit{D} et \textit{E} %
extremitatibus %
chordae \textit{DACBE} sint affixa,%
\indexS{pondus affixum chordae}
unum \textit{C} medio alicui puncto chordae,%
\indexS{chorda tensa}
quaeritur situs punctorum talis
ut centrum commune%
\indexS{centrum gravitatis commune}
\edtext{%
gravitatis \textit{G} descendat%
\indexS{descensus centri gravitatis}
quantum potest.%
}{%
\lemma{gravitatis \textit{G}}%
\Bfootnote{%
\textit{(1)}~sit omnium %
\textit{(2)}~descendat quantum potest.%
~\textit{L}%
}}
Jungantur \textit{D} et \textit{E},
sitque centrum eorum \textit{F},
rursusque jungantur \textit{C} et \textit{F},
sitque centrum \textit{G}.%
\indexS{centrum gravitatis ponderum suspensorum}
Sint \textit{A} et \textit{B} in eadem horizontali,
id enim arbitrarium est%
\bracks{;}
\ledpb%
et ipsi \textit{AB} parallela seu horizontalis
transeat per \textit{G},
occurrens
ipsi \textit{AD} %
in \textit{H}.
\edlabel{LH_35_11_14_015_015v_evkq-1}%
Quaeritur situs punctorum
\textit{D}, \textit{C}, \textit{E} talis %
ut \textit{AH} %
sit omnium possibilium maxima.%
\edlabel{LH_35_11_14_015_015v_evkq-2}%
\pend%
\begin{edfigure}[1.50em][0.00em]
\hspace{-65mm}%
\includegraphics[scale=1.0]%
{gesamttex/edit_VIII-4/images/LH_35_11_14_015_d1_015v.pdf}%
\edcaption{%
\hspace{-65mm}%
\bracks{\textit{Fig.\:1}}%
}%
\vspace{-12.25em}%
\hspace{64mm}%
\includegraphics[scale=1.0]%
{gesamttex/edit_VIII-4/images/LH_35_11_14_015_d2_015v.pdf}
\edcaption[0.75em]{%
\hspace{64mm}%
\bracks{\textit{Fig.\:2}}%
}%
\vspace{1.50em}%
\hspace{-65mm}%
\includegraphics[scale=1.0]%
{gesamttex/edit_VIII-4/images/LH_35_11_14_015_d3a_015v.pdf}%
\edcaption[0.0em]{%
\hspace{-65mm}%
\bracks{\textit{Fig.\:3a}}%
\pagelabel{LH_35_11_14_015v_Fig.3}%
\edtext{}{%
\lemma{\textit{Zum Diagramm} \bracks{\textit{Fig.\:3}}
\textit{(die Bezeichnung F wird vom Hrsg. durch $[F']$ ersetzt):}%
}%
\Afootnote[nonum]{%
\newline%
$\hash$
$\left.%
\protect\begin{array}{r}%
Q[F']X + X[F']V = \text{rect.}\\
R[F']V + V[F']X = \text{rect.}\protect\end{array}%
\right\}$
\protect\begin{tabular}{l}%
Ergo ang.
$Q[F']X = \text{ang.}\ 
R[F']V$%
\\
Ergo $\bigtriangledown$\textsuperscript{la} haec
$[F']RV$, $[F']QX$ similia%
\protect\end{tabular}%
\\
Jam vis
qua tenditur filum \textit{BC}%
\indexS{filum tensum}%
\indexS{vis qua tenditur filum}%
\indexS{tensio fili}
est ad vim
qua tenditur filum \textit{AC}
ut $[F']Q$ ad $[F']R$
seu ut $[F']X$ ad $[F']V]$.
\bracks{\textit{Gestrichen:}}
{\footnotesize%
$X[F']C = CBE = AC[F']$,
$V[F']C = CAD = AC[F']$,
$ACB = X[F']V$,
$XV = C[F']$
}%
}%
\lemma{\bracks{\textit{Fig.\:3}}}%
\Cfootnote[nonum]{%
Die im Diagramm und Text zweifach verwendete Bezeichnung \textit{F} wird vom Hrsg. desambiguiert.
In der ersten Fassung des Diagramms war der Punkt \textit{B} höher als der Punkt \textit{A} angesetzt.%
}%
}%
}%
\vspace{-13.6em}%
\hspace{66mm}%
\includegraphics[scale=1.00]%
{gesamttex/edit_VIII-4/images/LH_35_11_14_015_d3b_015v.pdf}%
\edcaption[0.5em]{%
\hspace{66mm}%
\bracks{\textit{Fig.\:3b}}%
}
\end{edfigure}
\newpage%
\pstart%
\edlabel{LH_35_11_14_015_015v_rnd-1}%
\edtext{}{%
\xxref{LH_35_11_14_015_015v_rnd-1}{LH_35_11_14_015_015v_rnd-2}%
\lemma{\textit{Am Rand:}}%
\Afootnote{%
\textit{d} hic mihi non differentiam%
\indexS{differentia}
sed quantitatem%
\indexS{quantitas}
significat.}}%
Pondera \textit{C}, \textit{D}, \textit{E} %
designentur per literas \textit{c}, \textit{d}, \textit{e},%
\indexS{litera designans pondus}
et
\edtext{%
\textit{AB}, \textit{h},%
}{%
\lemma{\textit{AB}, \textit{h},}%
\Bfootnote{%
\textit{erg.\;L}%
}}
\textit{AD} vocetur \textit{a}
et \textit{BE},\;\textit{b}.
Et transeunte \textit{LCM} horizontali per \textit{C},
tunc \textit{AL} vocetur \textit{l}
et \textit{BM},
\textit{m}.
Datur et \textit{DAC},
quae vocetur \textit{q},
et \textit{EBC},
quae vocetur \textit{r}.
Erit
$ CA = q-a $,
et
$ CB = r-b $,
et fiet
$ CL = \sqrt{qq-2qa+aa-ll} $
et
$ CM = \sqrt{rr-2rb+bb-mm} $.
Ex \textit{E} et \textit{F} demittantur %
perpendiculares \textit{EN} et \textit{FP} in \textit{AD}.
Est autem \textit{EN}, seu \textit{AB}, \textit{h};
erit \textit{FP} ad \textit{NE},
seu ad \textit{h},
ut \textit{e} ad
$ \overline{d+e} $,
seu
$ FP = he : \overline{d+e} $;
et similiter %
\textit{DP} ad \textit{DN},
seu ad $ \overline{a-b} $,
ut \textit{e} ad $ \overline{d+e} $.
Ergo fit
$ DP =
\overline{a-b}\;e : \overline{d+e} $.
Similiter \textit{LH} ad \textit{LP} %
ut \textit{CG} ad \textit{CF},
seu ut $\overline{d+e}$ ad $\overline{c+d+e}$.
Est autem $LP=LD-DP$ et $LD=a-l$,
et \textit{DP} inventa est.
Ergo
fit
$ LP = a - l - \overline{a-b}\, e : \overline{d+e} $. %
Ergo fit \textit{LH} ad \textit{LP},
seu ad
$ \overline{\overline{a-l}\;\overline{d+e} - \overline{a-b}\,e} : \overline{d+e} $,
ut
$ \overline{d+e} $ ad $ \overline{c+d+e} $.
Ergo
$ LH = \overline{\overline{a-l}\;\overline{d+e} - \overline{a-b}\,e} : \overline{c+d+e} $.
Cui si addas \textit{AL} seu \textit{l},
fiet \textit{AH} talis:
$ AH = cl\, %
\ovalbox{+\,dl+el}%
+ ad\, %
\ovalbox{+\,ae}%
\ovalbox{-\,dl-el}%
\ovalbox{-\,ae}%
+\, be, : \overline{c+d+e} $,
et fit
$ AH = cl + ad + be, : \overline{c+d+e} $.
Quod etiam praevideri potuerat.%
\edlabel{LH_35_11_14_015_015v_rnd-2}%
\pend%
\pstart%
\edtext{%
Quaeritur \textit{AH} omnium possibilium minima,%
\indexS{minima omnium possibilium}%
}{%
\lemma{Quaeritur \bracks{...} minima}%
\Cfootnote{%
Gesucht wurde der maximale Wert; %
vgl. \refpassage{LH_35_11_14_015_015v_evkq-1}{LH_35_11_14_015_015v_evkq-2}.%
}}
fit $ cl + ad + be, : \overline{c+d+e} = m $, 
\edtext{%
\edtext{%
ergo
$ c\,d\overline{l}+d\,d\overline{a}+e\,d\overline{b} = 0$.%
}{%
\lemma{ergo}%
\Bfootnote{%
\textit{(1)}~$c\,d\overline{l}+d\overline{d}+d\overline{b}=0$ %
\textit{(2)}~$c\,d\overline{l}+d\,d\overline{a}+e\,d\overline{b}=0$%
~\textit{L}%
}}%
}{%
\lemma{\textit{Zur Rechnung:}}%
\Afootnote{\ %
$dx^e=e\cdot x^{e-1}dx$\ \
\bracks{/}\ \
$d\sqrt[2]{x}=dx^{1:2}=\dfrac{1}{2}x^{-1:2}dx=\dfrac{dx}{2\sqrt{x}}$%
}}
(Res eo redit,
ut quaeratur uti planum
superficiem tangens
plano dato parallelam,%
\indexS{planum tangens superficiem parallelam}
tunc ordinata in superficie ad planum erit minima.)%
\indexS{ordinata minima}
Caeterum subest adhuc aliqua relatio,%
\indexS{relatio}
\edtext{%
nam sit
$ CL = v $,%
}{%
\lemma{nam}%
\Bfootnote{%
\textit{(1)}~\textit{AD} seu %
\textit{(2)}~$ d + e $ %
\textit{(3)}~sit $ CL = v $,%
~\textit{L}%
}}
fiet:
$ a + \sqrt{ll+vv} = q $,
et
$ b +
\edtext{%
\sqrt{h^2-2hl+ll+vv}%
}{%
\lemma{$\sqrt{h^2-2hl+ll+vv}$}%
\Cfootnote{%
Der Summand $-\,2hl$ heißt $-\,2hv$.
Der Fehler wirkt sich auf die %
Herleitung bis zu \refpassage{LH_35_11_14_015_015r_prdx}{LH_35_11_14_015_015r_prdx} aus.
Weitere Rechnungsfehler kommen dabei hinzu.%
}}
 =\,r $.
Hos valores substituamus in aeq.%
\indexS{valor substituendus}%
\indexS{aequatio}
$ cl + ad + be = m\,\overline{c+d+e} $,
\edtext{%
fiet
\edlabel{LH_35_11_14_015_aequpp-a}%
$ cl + d\,\overline{q-\sqrt{ll+vv}} + e\,\overline{r-\sqrt{hh-2hl+ll+vv}} = m\,\overline{c+d+e} $%
\edlabel{LH_35_11_14_015_aequpp-b}
\bracks{\textit{%
kleingedruckter Text 
gestr. u. abbrechend:}}
{\footnotesize{%
et fiet
$ c - c - dl :
\edtext{%
\sqrt{ll+vv} + l\;\overline{h-l} : \sqrt{hh-2hl+ll+vv} = 0 $
seu
$ c = d\,\overline{AL:LC} - e\,CM : $%
}{%
\lemma{$\sqrt{ll+vv}$}%
\Bfootnote{%
\textit{(a)}~$ -\, l +\, h $ %
\textit{(b)}~$ -\, \overline{l-h} $ %
\textit{(c)}~$ +\, l\;\overline{h-l} : \sqrt{hh-2hl+ll+vv} = 0 $ seu $ c = $ %
\textit{(aa)}~\textit{c} %
\textit{(bb)}~$d\,\overline{AL:LC}-e\,CM:$%
~\textit{L}%
}}%
}}%
}{%
\lemma{\textit{Zur Rechnung, zwischenzeilig:}}%
\Afootnote{%
Jam error%
\indexS{error}%
\bracks{:}
non est $=\,c$.
Pro \textit{h}\angles{\textit{l}} pone \textit{h}\angles{\textit{v}}.%
}}%
\pend%
\pstart%
Caeterum
quia possumus \textit{a} et \textit{b} assumere pro arbitrio,%
\indexS{arbitrium}
ideo cum \textit{q} et \textit{r} sint datae,
sumamus ita
ut $q-a$ sit $=r-b$,
seu ut fiat $AC=BC$;
et fiet
$ll+vv=hh-2hl+ll+vv$
seu
\rule[-0mm]{0pt}{4mm}%
$l=\dfrac{1}{2}h$,
et fiet
$ch:2+d\,\overline{q-\omega}+e\,\overline{r-\omega}=m\,\overline{c+d+e}$,
posito $ \omega = AC = BC $
et fit $ -\, d - e = 0 $.
Eo casu debet esse
$ d = e $.
\bracks{%
\href{\Scan{LH_35_11_14_015r}}%
{15\,r\textsuperscript{o}}%
}
\pend%
\newpage%
\pstart%
Maxima ex infinitis maximis hic quaerenda.%
\indexS{maxima quaerenda ex infinitis maximis}
\pend%
\pstart%
\edlabel{LH_35_11_14_015_015r_fatio-1}%
\edtext{%
Facius Dullerius%
\indexN{\textls{Fatio de Duillier} (Facius Dullerius), Nicolas 1664\textendash1753}
casum specialissimum tantum,%
\indexS{casus specialissimus}
et hunc recte resolvit.%
}{%
\lemma{Facius \bracks{...} resolvit}%
\Cfootnote{%
Siehe zu dieser Anspielung die ausführlichen Datierungsgründe unten, S.\;\pageref{58943_dat}\,f.
}}%
\edlabel{LH_35_11_14_015_015r_fatio-2}
Possumus
loco infinitorum circuituum%
\indexS{circuitus infiniti}
concipere
\edtext{%
ponder\bracks{a}%
}{%
\lemma{pondere}%
\Bfootnote{%
\textit{L\;ändert Hrsg.}%
}}
in quacunque proportione,
et flatum%
\indexS{flatus sustinens pondus}
vel jactum aquae%
\indexS{jactus aquae sustinens pondus}
ponder\angles{a}
sustinentem,%
\indexS{pondus sustentum flatu vel jactu aquae}
linea jactus est perpendicularis ad curvam,
vis jactus autem%
\indexS{vis jactus aquae}
mensuranda per centrum gravitatis.%
\indexS{vis mensuranda per centrum gravitatis}
Idem est
si ad potentias ascendatur.%
\indexS{potentia}%
\indexS{ascensus ad potentias}%
\edlabel{LH_35_11_14_015_015r_fatio-3}
Probatur
quod corpus trahitur ea vi,%
\indexS{vis trahens corpus}
quia,
qua trahit trahitur.%
\indexS{corpus tractum vi}
Aequali autem vi retinacula trahit,%
\indexS{retinaculum}%
\indexS{corpus trahens retinacula}
nec longitudo chordae%
\indexS{chorda tracta corpore}
ad rem
\edlabel{35_11_14_015_2a}%
facit.%
\edtext{}{%
\xxref{35_11_14_015_2a}{35_11_14_015_2b}%
\lemma{facit.}%
\Bfootnote{%
\bracks{/}
\textit{(1)}~Si fila con
\textit{(2)}~Si catenae vel fila constent%
~\textit{L}%
}}
\pend%
\pstart%
Si catenae%
\indexS{catena constans ex articulis tensis}
vel fila%
\indexS{filum constans ex articulis tensis}%
\indexS{filum tensum}
constent%
\edlabel{35_11_14_015_2b}
ex articulis rectilineis
iisque tensis,%
\indexS{articulus tensus catenae}\indexS{articulus tensus fili}
stylo per fila incedente%
\indexS{stylus incedens per filum}
describatur
\edtext{%
curva%
\indexS{curva descripta a stylo}
sitque ea articulorum dispositio,%
\indexS{dispositio articulorum in filo}
ut omnes lineam rectam faciant
excepto articulo fili%
\indexS{articulus tensus catenae}\indexS{articulus tensus fili}
(unius replicationis)%
\indexS{replicatio fili}
qu\bracks{i}
a stylo attingitur,%
\indexS{stylus incedens per filum}
utique patet eum%
}{%
\lemma{curva}%
\Bfootnote{%
\textit{(1)}~constat %
\textit{(2)}~tum manifestum est ut funiculum quem stylus %
\textit{(a)}~attingit %
\textit{(b)}~solum attingit %
\textit{(3)}~sitque ea %
\bracks{...} (unius replicationis) %
\bars{quae \textit{ändert Hrsg.}} a stylo attingitur, utique %
\textit{(a)}~is articulus %
\textit{(b)}~patet eum%
~\textit{L}%
}}
continuatum esse tangentem curvae.%
\indexS{tangens curvae}%
\indexS{curva descripta a stylo}
Porro idem est
si plura sint fila%
\indexS{filum replicatum}
\edtext{%
replicata,
articuli enim plurium filorum%
\indexS{articulus tensus catenae}\indexS{articulus tensus fili}%
}{%
\lemma{replicata,}%
\Bfootnote{%
\textit{(1)}~duo enim articuli %
\textit{(2)}~articuli enim plurium filorum%
~\textit{L}%
}}
in puncto tensionis%
\indexS{punctum tensionis in filo}%
\indexS{tensio fili}
seu concursus eorum inter se%
\indexS{punctum concursus articulorum in filo}
concipi possunt congruentes,
cum nihil referat
quae sit eorum
\edlabel{LH_35_11_14_015_015r_khj-1}%
magnitudo.%
\edtext{}{%
\xxref{LH_35_11_14_015_015r_khj-1}{LH_35_11_14_015_015r_khj-2}%
\lemma{magnitudo.}%
\Bfootnote{%
\bracks{/}
\textit{(1)}~Intelligi a %
\textit{(2)}~Porro%
~\textit{L}%
}}%
\pend%
\begin{edfigure}[2.0em][0.00em]
\includegraphics[scale=1.0]%
{gesamttex/edit_VIII-4/images/LH_35_11_14_015_d4_015r.pdf}
\edcaption[0.25em]{%
\bracks{\textit{Fig.\:4}}%
\edtext{}{%
\lemma{\textit{Zum Diagramm} \bracks{\textit{Fig.\:4}}:}%
\Afootnote[nonum]{%
Quid si foci focorum?%
\indexS{focus focorum}%
\indexS{focus curvae}
ut%
\textsuperscript{\bracks{a}}
fune \textit{ACDB} teneatur recta rigida \textit{CD}
sic ut per funem currere possit.%
\indexS{funis tensus}
Ipsa autem recta \textit{CD} novos ferat focos
ut \textit{E},\;\textit{F}
quibus alligatus funis \textit{EGF} stylo \textit{G} tendatur et \textit{GG} sit curva.
\\[1.5ex]%
{\footnotesize%
\textsuperscript{\bracks{a}}~ut %
\textit{(1)}~funibus %
\textit{(2)}~fune %
\textit{(a)}~\textit{ACBD} %
\textit{(b)}~\textit{ACDB}%
~\textit{L}%
}%
}}%
}%
\end{edfigure}
\newpage%
\pstart%
\edlabel{LH_35_11_14_015_015r_bve-1}%
Porro%
\edlabel{LH_35_11_14_015_015r_khj-2}
cum stylus nitens%
\indexS{stylus nitens in filum}
seu filum intendens%
\indexS{stylus intendens filum}%
\indexS{filum tensum}
unum tantum articulum attingit,%
\indexS{articulus tensus catenae}\indexS{articulus tensus fili}
est ad ipsum perpendicularis.
Alioqui enim
si obliquus esset
\edtext{%
articulus \textit{ar}%
\indexS{articulus tensus catenae}\indexS{articulus tensus fili}
ad stylum \bracks{\textit{st}} nitentem%
\indexS{stylus nitens in filum}%
}{%
\lemma{articulus}%
\Bfootnote{%
\textit{(1)}~, in eo stylus %
\textit{(2)}~\textit{ar} ad stylum %
\bars{\textit{s} \textit{ändert Hrsg.}} nitentem%
~\textit{L}%
}}
linea \textit{ns};
manifestum est
stylum posse adhuc progredi
versus angulum \textit{r}.
Nam si dicas
fieri posse
ut stylus \textit{st} totum angulum
\edtext{%
impleat,
respondeo detrahendam%
}{%
\lemma{impleat}%
\Bfootnote{%
\textit{(1)}~ajo detra %
\textit{(2)}~respondeo detrahendam%
~\textit{L}%
}}
esse magnitudinem inutilem,
cum assumi possit quantumvis parvus%
\indexS{angulus quantumvis parvus}
et velut in puncto consistens.
Perpendiculari ergo styli nisu existente \textit{ps},%
\indexS{nisus styli in filum}%
\indexS{stylus nitens in filum}
manifestum est
stylum intendentem fila%
\indexS{stylus intendens filum}
\textit{fa}, \textit{oa}, \textit{cr}, \textit{vr}
ab ipsis in tantum retrahi,%
\indexS{stylus retractus a filo tenso}%
\indexS{filum retrahens stylum}
in quantum ipsa trahit,%
\indexS{stylus trahens filum tensum}%
\indexS{filum tractum a stylo}
adeoque perinde
\edtext{%
esse
(qu\bracks{i}a \textit{ar} concipitur
indefinitae parvitatis)%
\indexS{parvitas indefinita}
ac si stylus \bracks{\textit{st}} %
impelleretur conatibus parallelis%
\indexS{conatus paralleli}%
\indexS{conatus impellens stylum}%
\indexS{stylus impulsus conatibus parallelis}
ipsis \textit{af}, \bracks{\textit{ao}}, \textit{rc}, \textit{rv}
ipsius styli conatum aequilibrantibus.%
\indexS{conatus styli intendentis filum}%
\indexS{conatus aequilibrati}%
}{%
\lemma{esse}%
\Bfootnote{%
\textit{(1)}~ac si co %
\textit{(2)}~( %
\bars{qua \textit{ändert Hrsg.}} \textit{ar} concipitur %
\bracks{...} si stylus %
\bars{\textit{s} \textit{ändert Hrsg.}} impelleretur conatibus parallelis ipsis \textit{af}, %
\bars{\textit{of} \textit{ändert Hrsg.}} \textit{rc}, \textit{rv} ipsius styli conatum aequilibrantibus.%
~\textit{L}%
}}
\pend%
\pstart%
Cum autem conatus isti sint aequales;%
\indexS{conatus aequales}%
\indexS{conatus impellens stylum}
hinc perinde est,
ac si centro \textit{s}, vel \textit{a},
\edtext{%
vel \textit{r}
(haec enim puncta coincidunt)
describatur%
}{%
\lemma{vel \textit{r}}%
\Bfootnote{%
\textit{(1)}~describat %
\textit{(2)}~(haec enim puncta coincidunt) describatur%
~\textit{L}%
}}
circulus secans rectas
\textit{af}, \textit{ao}, \textit{rc}, \textit{rv}
in punctis \textit{i},
\edtext{%
$[r']$,%
}{%
\lemma{\textit{r}}%
\Bfootnote{%
\textit{L\;ändert Hrsg.}%
}}
\textit{l}, \textit{m}
eruntque nisus ut \textit{ai},%
\indexS{nisus styli in filum}
\edtext{%
$a[r']$,%
}{%
\lemma{\textit{ar}}%
\Bfootnote{%
\textit{L\;ändert Hrsg.}%
}}
\textit{rl}, \textit{rm}
et nisus ipsis aequilibratus est,%
\indexS{nisus styli in filum}%
\indexS{nisus aequilibratus}
cujus directio est \textit{ps},%
\indexS{directio nisus}
transiens per centrum gravitatis%
\indexS{nisus transiens per centrum gravitatis}
omnium punctorum \textit{i},
\edtext{%
$[r']$,%
}{%
\lemma{\textit{r}}%
\Bfootnote{%
\textit{L\;ändert Hrsg.}%
}}
\textit{l}, \textit{m},%
\indexS{centrum gravitatis omnium punctorum}
\edtext{%
sed multiplus%
\indexS{nisus multiplus}%
}{%
\lemma{sed}%
\Bfootnote{%
\textit{(1)}~multiplicatus %
\textit{(2)}~multiplius %
\textit{(3)}~multiplus%
~\textit{L}%
}}
ipsius nisus \textit{ps}
in ratione numeri punctorum \textit{i},
\edtext{%
$[r']$,%
}{%
\lemma{\textit{r}}%
\Bfootnote{%
\textit{L\;ändert Hrsg.}%
}}
\textit{l}, \textit{m} ad unitatem.%
\edlabel{LH_35_11_14_015_015r_bve-2}
Idem est
si plura puncta,
ut \textit{v} et \textit{c},
coincidant inter se%
\bracks{;}
considerabuntur enim ponderis
\edtext{%
respectu velut%
}{%
\lemma{respectu}%
\Bfootnote{%
\textit{(1)}~qua %
\textit{(2)}~velut%
~\textit{L}%
}}
plura.
\pend%
\begin{edfigure}[2.00em][0.00em]
\hspace{-45mm}%
\includegraphics[scale=1.0]%
{gesamttex/edit_VIII-4/images/LH_35_11_14_015_d5a_015r.pdf}
\edcaption[0.0em]{%
\hspace{-45mm}%
\bracks{\textit{Fig.\:5a}}%
\edtext{}{%
\lemma{\bracks{\textit{Fig.\:5a}}}%
\Cfootnote[nonum]{%
Die im Diagramm und Text zweifach verwendete Bezeichnung \textit{r} wird vom Hrsg. desambiguiert.%
}}%
}%
\vspace{-11.250em}%
\hspace{45mm}%
\includegraphics[scale=0.71]%
{gesamttex/edit_VIII-4/images/LH_35_11_14_015_d5b_015r.pdf}
\edcaption[0.35em]{%
\hspace{45mm}%
\bracks{\textit{Fig.\:5b}}%
}%
\end{edfigure}%
\newpage%
\pstart%
Caeterum resumamus
\edtext{%
aequationem paginae praecedentis,%
\indexS{aequatio resumta}%
\indexS{pagina}%
}{%
\lemma{aequationem%
}%
\Cfootnote{%
Vgl. \refpassage{LH_35_11_14_015_aequpp-a}{LH_35_11_14_015_aequpp-b}.%
}}%
\edtext{ %
et pro %
}{%
\lemma{praecedentis}%
\Bfootnote{%
\textit{(1)}~quae erat \textit{c} %
\textit{(2)}~et pro%
~\textit{L}%
}}%
$ ll + vv $ scribamus
$ \omega\omega $,
et fiet %
$ cl + dq - d\omega + er - e \sqrt{hh-2hl+\omega\omega} \stackrel{\astrosun}{=} m\,\overline{c+d+e} $;
posita \textit{m}
\edtext{%
maxima%
\bracks{,}
fiet
$ c + eh%
}{%
\lemma{maxima; fiet}%
\Bfootnote{%
\textit{(1)}~$ c + h : $ %
\textit{(2)}~$ c + h\, - $ %
\textit{(3)}~$ c + eh : $%
~\textit{L}%
}}
: \sqrt{hh-2hl+\omega\omega} = 0 $,
seu
$ \sqrt{hh-2hl+\omega\omega} = -\,eh : c $.
Et fieret
$ \sqrt{hh-2hl+\omega\omega} $
seu \textit{BC}, ad \textit{h} seu \textit{AB},
ut \textit{e} ad
\edtext{%
\bracks{$-\,c$}.%
}{%
\lemma{$-\,c$}%
\Bfootnote{%
\textit{erg. Hrsg.}%
}%
}
Ergo pari jure erit
$ \sqrt{ll+vv} $
seu $ \omega $
seu \textit{AC} ad \textit{h} seu \textit{AB},
ut \textit{d} ad \textit{c}.
Quod
\edtext{%
debet
\bracks{oriri}
tollendo \textit{l} ex aequ. superiore%
\indexS{aequatio superior}
ope valoris inventi%
\indexS{valor inventus}%
}{%
\lemma{debet}%
\Bfootnote{%
\bars{oriri \textit{gestr., wieder gültig macht Hrsg.}} %
\textit{(1)}~substituendo valorem\indexS{valor substituendus} %
\textit{(2)}~tollendo %
\textit{(a)}~$ \omega $ %
\textit{(b)}~\textit{l} ex %
\bracks{...} valoris inventi%
~\textit{L}%
}}
et rursus quaerendo maximam,%
\indexS{maxima quaesita}
\edlabel{LH_35_11_14_015_015r_wlj-3a}%
\edtext{%
nempe fit aequ.%
\indexS{aequatio}
\edtext{%
$ hh - 2hl + \omega\omega = ee\,hh : cc. $%
}{%
\lemma{\textit{Am Rand:}}%
\Afootnote{%
Error in calculo,%
\indexS{error in calculo}%
\indexS{calculus}
quia pro 2\textit{hl} debuisset scribi 2\textit{hv}.%
}}
seu
$ l = h : 2 + \omega\omega : 2h - hee : 2cc $,%
\edlabel{LH_35_11_14_015_015r_wlj-3b}%
}{%
\lemma{nempe}%
\Bfootnote{%
\textit{(1)}~$ \omega = \sqrt{\llcorner e^2h^2:cc-hh+2hl\lrcorner} $ %
\textit{(2)}~fit aequ. %
\textit{(a)}~$ c + e $ %
\textit{(b)}~$ cl + dq $ %
\textit{(c)}~$ hh - 2hl + \omega\omega = ee\,hh : cc $, %
\textit{(aa)}~seu %
\textit{(aaa)}~$ l = eehh : cc $ %
\textit{(bbb)}~$ \dfrac{hh}{2} + \dfrac{\omega\omega}{2h} - eeh : cc $ %
\textit{(bb)}~seu $ l = h : 2 + \omega\omega : 2h - hee : 2cc $.%
~\textit{L}%
}}
Unde ex aeq. \astrosun\ habemus:
$ ch : 2 + c\,\omega\omega : 2h - hee : 2c + dq - d \omega + er + eeh : c = m\cdot\overline{c+d+e} $
seu
$ ch : 2 + c\,\omega\omega : 2h + hee : 2c + dq - d \omega + er = m\,\overline{c+d+e} $,
et fiet
$ c\,\omega : h - d = 0 $,
seu
$\omega$ ad \textit{h} ut \textit{d} ad \textit{c}.
Ut praediximus.%
\edlabel{LH_35_11_14_015_015r_prdx}%
\pend%
\pstart%
Et notanda est haec methodus
ex infinitis maximis rursus maximam%
\indexS{methodus quaerendi maximam ex infinitis maximis}%
\indexS{maxima quaerenda ex infinitis maximis}
\edtext{%
quaerendi.
Quando scilicet
\bracks{sunt}
plures indeterminata\bracks{e}
quae ex sola conditione maximae quaerendae%
\indexS{conditio maximae quaerendae}
determinari debent,
ut hoc loco \textit{l} et $\omega $,
tunc omnes%
}{%
\lemma{quaerendi.}%
\Bfootnote{%
\textit{(1)}~Unam %
\textit{(2)}~Omnes scilicet %
\textit{(3)}~Quando scilicet %
\bars{sunt \textit{erg. Hrsg.}} plures %
\bars{indeterminatas \textit{ändert Hrsg.}} quae ex %
\textit{(a)}~solo %
\textit{(b)}~sola data %
\textit{(c)}~sola conditione \bracks{...} tunc omnes%
~\textit{L}%
}}
habeantur pro determinatis
demta una,
quae exposita ad maximam aequatione inveniatur.%
\indexS{aequatio exposita ad maximam}
Et reperta per aequationem incrementalem%
\indexS{aequatio incrementalis}
tollatur ex aequatione.
Ex residuis duabus rursus
una tollatur eadem methodo
et ita porro%
\bracks{;}
vel quod eodem redit%
\bracks{,}
statim prima aequatio
secundum quamlibet indeterminatarum%
\indexS{indeterminata in aequatione}
quasi solam talem
tractetur.
\edtext{%
Vera igitur methodus%
\indexS{methodus vera}
sed error in calculo.%
\indexS{calculus}%
\indexS{error in calculo}
Itaque resumamus:%
}{%
\lemma{Vera igitur \bracks{...} Itaque resumamus:}%
\Bfootnote{%
\textit{erg.\;L}%
}}
\pend%
\pstart%
\edtext{}{%
\lemma{\textit{Am Rand:}}%
\Afootnote{%
Quae sequuntur sufficiunt
et errore calculi carent.%
\indexS{calculus}%
\indexS{error calculi}%
}}%
\edtext{%
\edtext{%
Pondera \textit{c}, \textit{d}, \textit{e}.%
\indexS{pondus suspensum}%
}{%
\lemma{Pondera \textit{c}, \textit{d}, \textit{e}}%
\Cfootnote{%
Vgl. das Diagramm \bracks{\textit{Fig.\:6}}, das an %
\bracks{\textit{Fig.\:3}} anknüpft.%
}}
Eorum descensus%
\indexS{descensus ponderis}
\textit{AL}, \textit{l}.%
}{%
\lemma{Pondera}%
\Bfootnote{%
\textit{(1)}~\textit{d}, \textit{e} %
\textit{(2)}~\textit{c}, \textit{d}, \textit{e}. %
\textit{(a)}~Pond %
\textit{(b)}~Eorum descensus %
\textit{(aa)}~\textit{AD} %
\textit{(bb)}~\textit{AL} \textit{l}.%
~\textit{L}%
}}
\textit{AD}, \textit{a}.
\textit{BE}, \textit{b}.
\textit{AF}, \textit{f}.
\edtext{%
\textit{AG}, \textit{m}.%
}{%
\lemma{\textit{AG},}%
\Bfootnote{%
\textit{m}.
\textit{erg.\;L}%
}}
\textit{BF}, \textit{h}.
$ \overline{DA + AC} $, \textit{q}.
$ \overline{EB + BC} $, \textit{r}.
\textit{CL}, \textit{v}.
\textit{AC}, \textit{x}.
\textit{BC}, \textit{y}.%
\rule[0.0mm]{0.0mm}{4.6mm}
Ergo \textit{AG}
descensus centri gravitatis%
\indexS{descensus centri gravitatis}
ponderum%
\indexS{centrum gravitatis ponderum suspensorum}
erit
$ lc + ad + be , : \overline{c+d+e} \stackrel{(1)}{=} m $
\edtext{%
\bracks{qui}%
}{%
\lemma{quia}%
\Bfootnote{%
\textit{L\;ändert Hrsg.}%
}}
debet esse omnium possibilium maximus.%
\indexS{descensus maximus omnium possibilium}
Jam considerandum
\edtext{%
est
longitudinem funis \textit{DAC},
\ledpb%
item \bracks{\textit{CBE}}
esse datas.%
}{%
\lemma{est}%
\Bfootnote{%
\textit{(1)}~esse %
\textit{(2)}~longitudinem funis \textit{DAC}, item %
\bars{\textit{CDE} \textit{ändert Hrsg.}}, esse datas.%
~\textit{L}%
}}
Sit $ DA + AC = q $
et $ EB + BC = r $.
Jam \textit{AD} est \textit{a}
et $ AC^2$ est $AL^2+CL^2 $.
Jam \textit{CL} sit \textit{v},
fiet $ AC = \sqrt{ll+vv} = x $
et habemus
$ a + \sqrt{ll+vv} \stackrel{(2)}{=} q $.
Rursus \textit{BE} est \textit{b}
et $ BC^2 $ est $ BM^2 + CM^2 $.
Jam $ BM = AL - AF $
seu $ l - f $
et $ CM = FB - LC = h - v $,
et fiet
 $ b + \sqrt{ll-2lf+f\!f+hh-2hv+vv} \stackrel{(3)}{=} r $.
Ergo pro \textit{a} et \textit{b}
valores ex aeq.
\edtext{%
\bracks{2} et \bracks{3}%
}{%
\lemma{1 et~2}%
\Bfootnote{%
\textit{L\;ändert Hrsg.}%
}}
substituendo in aeq.\;1,%
\indexS{valor substituendus}%
\indexS{aequatio}
fiet
$ lc + dq -
\underbrace{d\sqrt{ll+vv}}_{\displaystyle\text{seu \textit{dx}}}
\edtext{%
+ ef%
}{%
\lemma{$+\, ef$}%
\Cfootnote{%
Der Summand ist falsch hergeleitet.
Der Fehler wirkt sich auf die folgende Herleitung
bis zu \refpassage{LH_35_11_14_015_015r_fhl}{LH_35_11_14_015_015r_fhl} aus.%
}}
+ er - \underbrace{e\, \sqrt{ll-2lf+f\!f+hh-2hv+vv}}_{\displaystyle\text{seu}\ ey} \stackrel{(4)}{=} m\,\overline{c+d+e} \stackrel{(5)}{=} lc + dq - dx + ef + er - ey $.
Ergo ponendo primum \textit{v} constantem,
fit
$ c - dl : x - el : y + ef : y \stackrel{(6)}{=} 0 $.%
\edlabel{LH_35_11_14_015_015r_fhl}
At ponendo \textit{l} constantem,
fit:
$ -\,dv : x + eh : y - ev : y%
\stackrel{(7)}{=} 0 $.
Ex aeq.\;7 %
patet
\edtext{%
esse
$ dv : e\,\overline{h-v} \: \squaredots \: x : y$,%
}{%
\lemma{esse}%
\Bfootnote{%
\textit{(1)}~$ l : e \: \squaredots \: \overline{h-v} $ seu $ CM : CB $ %
\textit{(2)}~$ dv : e\,\overline{h-v} \: \squaredots \: x : y $,%
~\textit{L}%
}}
seu momenta %
ponderum \textit{D} et \textit{E},%
\indexS{momentum ponderis suspensi}
ex axe \textit{HC}%
\indexS{axis}
\edtext{%
vel ex centro \textit{H},%
}{%
\lemma{vel}%
\Bfootnote{%
ex centro \textit{H} %
\textit{erg.\;L}%
}}
sunt ut chordae \textit{AC}, \textit{BC}.%
\indexS{chorda tensa}
In casu igitur chordarum aequalium,
pondera \textit{D} et \textit{E}
ex \textit{H}
forent in aequilibrio.%
\indexS{aequilibrium ponderum suspensorum}
\pend%
\pstart%
Quod an receptae propositioni consentiat%
\indexS{propositio recepta}
\edtext{%
videamus.
\edlabel{LH_35_11_14_015_015r_nempe-1}%
\edtext{}{%
\lemma{Nempe \bracks{...} aequalium}%
\xxref{LH_35_11_14_015_015r_nempe-1}{LH_35_11_14_015_015r_nempe-2}%
\Cfootnote{%
Die Ausführung knüpft letztlich an das Diagramm \bracks{\textit{Fig.\:3}} (S.\;\pageref{LH_35_11_14_015v_Fig.3}) an, wo die Bezeichnung \textit{F} zweifach verwendet ist.
Der Hrsg. desambiguiert durch Einsetzung der Bezeichnung $[F']$ anstelle von \textit{F}.%
}}%
Nempe%
}{%
\lemma{videamus.}%
\Bfootnote{%
\textit{(1)}~Pondus \textit{D} %
\textit{(2)}~Nempe%
~\textit{L}%
}}
si
\edtext{%
$ AC = BC $,%
}{%
\lemma{$ AC = BC $}%
\Cfootnote{%
Vgl. das Diagramm \bracks{\textit{Fig.\:7}},
das an %
\bracks{\textit{Fig.\:6}} und \bracks{\textit{Fig.\:3}} anknüpft.%
}}
debet
ex nostro hoc theoremate%
\indexS{theorema nostrum}
esse pondus \textit{D}
ad pondus \textit{E}%
\indexS{pondus suspensum}
ut \textit{BT} ad \textit{AS};
seu
\edtext{%
ex aliorum theoremate%
\indexS{theorema aliorum}%
}{%
\lemma{ex aliorum theoremate}%
\Cfootnote{%
Siehe \textsc{P.\;de Varignon}, \textit{Projet d'une nouvelle méchanique},
Proposition fondamentale des poids suspendus avec des cordes
(Paris 1687, S.\;11).\cite{02028}
Siehe zudem N.\;\ref{60365} im vorliegenden Band (Leibnizens Auszüge mit Bemerkungen aus Varignons \textit{Projet}).%
}}
est pondus \textit{D}
ad pondus \textit{E}%
\indexS{pondus suspensum}
ut \textit{\textup{[}F'\textup{]}R} ad \textit{\textup{[}F'\textup{]}Q}
\edtext{%
(posito angulos \textit{Q} et \textit{R} esse rectos),%
\indexS{angulus rectus}%
}{%
\lemma{(posito}%
\Bfootnote{%
angulos %
\bracks{...} esse rectos) %
\textit{erg.\;L}%
}}
ut
\edtext{%
ostendimus pagina praecedente%
\indexS{pagina}
ad signum\:$\hash$.%
\indexS{signum}%
}{%
\lemma{ ostendimus \bracks{...} signum\:$\hash$\,}%
\Cfootnote{%
Siehe die Randbemerkung zum Diagramm \bracks{\textit{Fig.\;3}},
S.\;\pageref{LH_35_11_14_015v_Fig.3}.
}}
Ergo debet esse \textit{BT} ad \textit{AS},
ut \textit{\textup{[}F'\textup{]}R} ad \textit{\textup{[}F'\textup{]}Q}.
Sed hoc reperitur verum.
Nam triangulum \textit{C\textup{[}F'\textup{]}A} est ad triangulum \textit{C\textup{[}F'\textup{]}B}
ut \textit{AS} ad \textit{BT}
(altitudines, ob communem basin \textit{C\textup{[}F'\textup{]}}).%
\indexS{altitudo trianguli}%
\indexS{basis trianguli}
Rursus triangulum \textit{C\textup{[}F'\textup{]}A} est ad triangulum \textit{C\textup{[}F'\textup{]}B}
ut \textit{\textup{[}F'\textup{]}Q} ad \textit{\textup{[}F'\textup{]}R}
(altitudines, ob aequales bases \textit{CA}, \textit{CB}).%
\indexS{altitudo trianguli}%
\indexS{basis trianguli}
Ergo \textit{AS} ad \textit{BT},
ut \textit{\textup{[}F'\textup{]}Q} ad \textit{\textup{[}F'\textup{]}R},
seu ut
$[F']X$ %
ad
$[F']V$ %
\edtext{%
(in fig.\;prima).%
\indexS{figura}%
}{%
\lemma{fig.\;prima}%
\Cfootnote{%
Das Diagramm \bracks{\textit{Fig.\:3}}.%
}}
\ledpb%
Sic ergo sunt pondera%
\indexS{pondus suspensum}
seu tensiones chordarum%
\indexS{tensio chordae}%
\indexS{chorda tensa}
in casu chordarum aequalium.%
\edlabel{LH_35_11_14_015_015r_nempe-2}
Ergo et semper,
quia
quae sit longitudo chordarum
nil refert.
\pend%
\pstart%
Igitur
quod peropportunum accidit%
\bracks{:}
ex sola aequatione~7
solvitur problema,%
\indexS{problema solutum}
nec opus est reliquis aequationibus~5,\;6 %
et praeterea explicatione ipsarum \textit{x} et \textit{y} uti.%
\indexS{explicatio}
Adeoque hoc
\edtext{%
modo semel tantum opus est
aequationis propositae~5%
\indexS{aequatio proposita}
depressione ad differentialem%
\indexS{aequatio differentialis}
seu incrementalem.%
\indexS{aequatio incrementalis}%
\indexS{depressio aequationis}%
}{%
\lemma{modo}%
\Bfootnote{%
\textit{(1)}~aequatione~5 differentiali %
\textit{(2)}~semel tantum \bracks{...} seu incrementalem.%
~\textit{L}%
}}
\pend%
\pstart%
Caeterum haec consentire
cum
\edtext{%
communi sententia%
\indexS{sententia communis}%
}{%
\lemma{communi sententia}%
\Cfootnote{%
Nicht ermittelt.%
}}
sic patet.
\edlabel{LH_35_11_14_015_015r_ost-1}%
\edtext{}{%
\xxref{LH_35_11_14_015_015r_ost-1}{LH_35_11_14_015_015r_ost-2}%
\lemma{Ostendimus \bracks{...} \textit{CX}, \textit{CV}, $C[F']$}%
\Cfootnote{%
Die Ausführung knüpft an das Diagramm \bracks{\textit{Fig.\:3}} (S.\;\pageref{LH_35_11_14_015v_Fig.3}) an, wo die Bezeichnung \textit{F} zweifach verwendet ist.
Der Hrsg. desambiguiert durch Einsetzung der Bezeichnung $[F']$ anstelle von \textit{F}.%
}}%
Ostendimus
\edtext{%
esse tensiones chordarum%
\indexS{tensio chordae}%
\indexS{chorda tensa}
\textit{AC} ad \textit{BC},
seu pondera \textit{D} ad \textit{E},%
\indexS{pondus suspensum}
ut $[F']R$ ad $[F']Q$,%
}{%
\lemma{esse}%
\Bfootnote{%
\textit{(1)}~pondera %
\textit{(2)}~tensiones chordarum \textit{AC} ad \textit{BC}, seu %
\textit{(a)}~chordas %
\textit{(b)}~pondera \textit{D} ad \textit{E}, %
\textit{(aa)}~ut $[F']Q$ %
\textit{(bb)}~ut $[F']R$ ad $[F']Q$,%
~\textit{L}%
}}
tum ex nostro calculo,%
\indexS{calculus noster}
tum etiam ex eo
\edtext{%
quia \textit{AB}
tracta a \textit{C}%
}{%
\lemma{quia}%
\Bfootnote{%
\textit{(1)}~rect
\textit{(2)}~libra
\textit{(3)}~\textit{AB} tracta a %
\bars{puncto \textit{gestr.}} \textit{C}%
~\textit{L}%
}}
per chordas \textit{AC}, \textit{BC}%
\indexS{chorda tensa}
est in
\edtext{%
aequilibrio,%
\indexS{aequilibrium ponderum suspensorum}
ergo%
}{%
\lemma{aequilibrio,}%
\Bfootnote{%
\textit{(1)}~quaerendo %
\textit{(2)}~ergo%
~\textit{L}%
}}
tractiones%
\indexS{tractio chordae}
\edtext{%
per \textit{CQ} et \textit{CR}
ut $[F']R$ ad $[F']Q$,%
}{%
\lemma{per}%
\Bfootnote{%
\textit{(1)}~\textit{CR} et \textit{CQ} %
\textit{(2)}~\textit{CQ} et \textit{CR} %
\textit{(a)}~reciproce %
\textit{(b)}~ut $[F']R$ ad $[F']Q$,%
~\textit{L}%
}}
seu ut $[F']V$ ad
\edtext{%
$[F']X$,
seu ut \textit{CX} ad \textit{CV}.
Quas rationes aequales esse%
\indexS{rationes aequales}
sic ostendo%
\bracks{:}%
}{%
\lemma{$[F']X$}%
\Bfootnote{%
\bars{seu ut \textit{CX} ad \textit{CV} \textit{erg.}}.
\textit{(1)}~Quod sic ostendo
\textit{(2)}~Quas rationes \bracks{...} sic ostendo%
~\textit{L}%
}}
angulus $Q[F']V$ est rectus%
\indexS{angulus rectus}
(quia rectus est ad \textit{Q},
et $[F']V$ parallela \textit{QC})
et similiter $X[F']R$ est rectus
(ob eandem rationem\bracks{)}.
Ergo anguli $Q[F']V$ et $X[F']R$ sunt aequales,
auferatur ab utroque communis 
\edtext{%
$[XF'V]$%
}{%
\lemma{$Q[F']R$}%
\Bfootnote{%
\textit{L\;ändert Hrsg.}%
}}
ergo et residui $X[F']Q$ et $V[F']R$ erunt aequales,
ergo triangula rectangula $[F']QX$ et $[F']RV$ sunt similia,%
\indexS{triangulum rectangulum}%
\indexS{triangula similia}
adeoque $[F']R$ ad $[F']Q$ ut $[F']V$ ad $[F']X$
\edtext{%
seu ut \textit{CX} ad \textit{CV},%
}{%
\lemma{seu}%
\Bfootnote{%
ut \textit{CX} ad \textit{CV}
\textit{erg.\;L}%
}}
adeoque tensiones funium \textit{AC}, \textit{BC},%
\indexS{tensio funis}%
\indexS{funis tensus}
seu pondera \textit{D}, \textit{E},%
\indexS{pondus suspensum}
\edtext{%
sunt [ut] $[F']R$ ad $[F']Q$,%
}{%
\lemma{sunt}%
\Bfootnote{%
\bars{ut \textit{erg. Hrsg.}} %
\textit{(1)}~\textit{HV} ad \textit{H} %
\textit{(2)}~$[F']R$ ad $[F']Q$,%
~\textit{L}%
}}
seu ut \textit{CX} ad \textit{CV}.
\pend%
\pstart%
Idem et alio adhuc%
\bracks{,}
nempe tertio modo%
\bracks{,}
sic probo breviter.
Pondus \textit{C}%
\indexS{pondus suspensum}
vi \textit{D}%
\indexS{vis ponderis suspensi}
tendit ad \textit{A}
et vi \textit{E} ad \textit{B},
tendit autem revera vi $[F']$,
ergo
\edtext{%
vires directrices%
\indexS{vis directrix}
ad \textit{A},
ad \textit{B},
ad $[F']$,
seu pondera
\textit{D}, \textit{E}, \textit{C}
sunt
(completo parallelogrammo $CX[F']V$)%
\indexS{parallelogrammum}
ut \textit{CX}, \textit{CV}, $C[F']$.%
}{%
\lemma{vires}%
\Bfootnote{%
\textit{(1)}~ad \textit{A} %
\textit{(2)}~directrices ad \textit{A}, ad \textit{B}, ad $[F']$, %
\textit{(a)}~sunt (completo parallelogrammo %
\textit{(b)}~seu pondera \textit{D}, \textit{E}, \textit{C} sunt %
\textit{(aa)}~ut %
\textit{(bb)}~(completo parallelogrammo $CX[F']V$) ut \textit{CX}, \textit{CV}, $C[F']$.%
~\textit{L}%
}}%
\edlabel{LH_35_11_14_015_015r_ost-2}%
\pend%
\newpage%
\begin{edfigure}[0.00em][0.0em]
\hspace{-53mm}%
\includegraphics[scale=1.0]%
{gesamttex/edit_VIII-4/images/LH_35_11_14_015_d6_015r.pdf}%
\edcaption[-0.50em]{%
\hspace{-53mm}%
\bracks{\textit{Fig.\:6}}%
}%
\vspace{-10.625em}%
\hspace{65mm}%
\includegraphics[scale=1.0]%
{gesamttex/edit_VIII-4/images/LH_35_11_14_015_d7_015r.pdf}%
\edcaption[1.25em]{%
\hspace{65mm}%
\bracks{\textit{Fig.\:7}}%
}%
\end{edfigure}
\vspace{4.5em}%
\pstart%
\noindent%
\centering%
*     *     *
\pend%
\vspace{5.0em}%
\selectlanguage{german}%
\frenchspacing%
\footnotesize%
\pstart%
\noindent%
\textls{Ausführliche Datierungsgründe:}%
\pagelabel{58943_dat}
Der Träger der undatierten Aufzeichnung N.\;\ref{58943} weist weder Spuren eines Wasserzeichens noch sonstige äußere datierungsrelevante Merkmale auf.
Im Text (\refpassage{LH_35_11_14_015_015r_fatio-1}{LH_35_11_14_015_015r_fatio-2}) spielt Leibniz allerdings -- in eher kryptischer Weise -- auf zwei Aufsätze des Schweizer Mathematikers Nicolas \textsc{Fatio de Duillier}\indexN{\textls{Fatio de Duillier} (Facius Dullerius), Nicolas 1664\textendash1753} an:
\textit{Réflexions sur une méthode de trouver les tangentes de certaines lignes courbes} (\textit{BUH}, Bd.\;V, April 1687, S.\,25--33)\cite{02539}\cite{02545} und
\textit{Réponse à} \angles{\textit{l'}}\textit{écrit de M.\;de T}\bracks{\textit{schirnhaus}} \textit{touchant une manière de déterminer les tangentes des lignes courbes} (\textit{BUH}, Bd.\;XIII, April 1689, S.\;46--76).\cite{02540}\cite{02545}
In den \textit{Réflexions}\cite{02539} wird die Methode zur Bestimmung der Tangenten, die Ehrenfried Walther von \textsc{Tschirnhaus}\indexN{\textls{Tschirnhaus} (H. v. Tch., H. v. Tsch.), Ehrenfried Walther v. 1651\textendash1708} in seiner Abhandlung \textit{Medicina mentis} (pars II, sectio II, regula III, Erläuterung zum zweiten Teil: Amsterdam 1687 \bracks{und 1686}, S.\;71--75)\cite{02542} dargelegt hatte, von Fatio de Duillier\indexN{\textls{Fatio de Duillier} (Facius Dullerius), Nicolas 1664\textendash1753} getadelt und verbessert.
In der \textit{Réponse}\cite{02540} erwidert der Schweizer Mathematiker ausführlich auf das Antwortschreiben, das Tschirnhaus\indexN{\textls{Tschirnhaus} (H. v. Tch., H. v. Tsch.), Ehrenfried Walther v. 1651\textendash1708} inzwischen, als Reaktion auf die \textit{Réflexions}, veröffentlicht hatte (\textit{BUH}, Bd.\;X, Juli 1688, S.\;497--509).\cite{02453}
Die von Fatio de Duillier\indexN{\textls{Fatio de Duillier} (Facius Dullerius), Nicolas 1664\textendash1753} kritisierte Methode Tschirnhausens\indexN{\textls{Tschirnhaus} (H. v. Tch., H. v. Tsch.), Ehrenfried Walther v. 1651\textendash1708} ging von Kurven aus, die mit einem Stift beschrieben werden können, der entlang eines an Brennpunkten gebundenen Fadens läuft (vgl. \textit{Medicina mentis}, pars II, sectio II, regula III, Erläuterung zum ersten Teil; S.\;68--71 samt Fig.\;5--11).\cite{02542} %
Bei seiner Anspielung in N.\;\ref{58943} dürfte sich Leibniz insbesondere auf die \textit{Réponse}\cite{02540} beziehen, wie seine Wortwahl in den unmittelbar folgenden Zeilen (\refpassage{LH_35_11_14_015_015r_fatio-2}{LH_35_11_14_015_015r_fatio-3}) zeigt:
seine Aussage \textit{Facius Dullerius casum specialissimum tantum, et hunc recte resolvit} entspricht Fatio de Duilliers\indexN{\textls{Fatio de Duillier} (Facius Dullerius), Nicolas 1664\textendash1753} Erläuterung \textit{j'ai cherché de propos déliberé un cas particulier, un exemple facile et de peu d'étendue}, welcher die Unzulänglichkeit von Tschirnhausens\indexN{\textls{Tschirnhaus} (H. v. Tch., H. v. Tsch.), Ehrenfried Walther v. 1651\textendash1708} Methode deutlich zeigen könne (\textit{Réponse}, S.\;48);\cite{02540}
im Ausdruck \textit{loco infinitorum circuituum} klingt Fatio de Duilliers\indexN{\textls{Fatio de Duillier} (Facius Dullerius), Nicolas 1664\textendash1753} Feststellung nach, gemäß Tschirnhausens\indexN{\textls{Tschirnhaus} (H. v. Tch., H. v. Tsch.), Ehrenfried Walther v. 1651\textendash1708} Methode erweise sich die Anwendung eines \textit{fil qui fit une infinité de retours} als unvermeidbar (ebd.);\cite{02540}
und Leibnizens Bemerkung \textit{Idem est si ad potentias ascendatur} knüpft an Fatio de Duilliers\indexN{\textls{Fatio de Duillier} (Facius Dullerius), Nicolas 1664\textendash1753} generalisierender Darstellung der eigenen Lösung an (vgl. ebd., Théoreme, S.\;63--66).\cite{02540}
Aus dieser Abhängigkeit von der \textit{Réponse}\cite{02540} ergibt sich der Spätfrühling 1689 als frühester möglicher Entstehungszeitpunkt der Aufzeichnung N.\;\ref{58943}.
\pend%
\pstart%
Wann genau Fatio de Duilliers\indexN{\textls{Fatio de Duillier} (Facius Dullerius), Nicolas 1664\textendash1753} Aufzätze \textit{Réflexions}\cite{02539} und \textit{Réponse}\cite{02540} von Leibniz gelesen wurden, ist allerdings nicht unmittelbar klar.
Dem Brief an Otto Mencke vom 22.\;Juli (1.\;August) 1687 lässt sich entnehmen, dass Leibniz zu jenem Zeitpunkt die \textit{Réflexions}\cite{02539} noch nicht kannte, obwohl er eine Ankündigung derselben in der \textit{Bibliothèque universelle et historique} (Bd.\;IV, Januar 1687, "Avertissement", S.\;\bracks{3\,r\textsuperscript{o}})\cite{02545} gelesen und diese sein Interesse geweckt hatte (\textit{LSB} I,\;4 N.\;541, S.\;645.12--14).\cite{02544}
Drei Monate später machte er sich aber auf den Weg nach Süddeutschland,\indexO{Süddeutschland} Österreich\indexO{Österreich} und Italien,\indexO{Italien} woher er erst Mitte Juni 1690 nach Hannover\indexO{Hannover} zurückkehrte.
Leibnizens erste briefliche Erwähnungen der Auseinandersetzung zwischen Fatio de Duillier\indexN{\textls{Fatio de Duillier} (Facius Dullerius), Nicolas 1664\textendash1753} und Tschirnhaus\indexN{\textls{Tschirnhaus} (H. v. Tch., H. v. Tsch.), Ehrenfried Walther v. 1651\textendash1708} stammen tatsächlich (nach heutigem Wissensstand) aus der Zeit nach der Italienreise:\indexO{Italien}
Im Schreiben an Jacob Bernoulli\indexN{\textls{Bernoulli}, Jacob 1655\textendash1705} vom 24.\;September~/ 4.\;Oktober 1690 hält er fest, %
anhand einer allgemeineren Methode ähnliche Ergebnisse gefunden zu haben wie die, die Fatio de Duilliers\indexN{\textls{Fatio de Duillier} (Facius Dullerius), Nicolas 1664\textendash1753} in seiner Kritik an Tschirnhaus\indexN{\textls{Tschirnhaus} (H. v. Tch., H. v. Tsch.), Ehrenfried Walther v. 1651\textendash1708} dargelegt habe (\textit{LSB} III,\;4 N.\;279, S.\;581.12--14).\cite{01260}
Und in seinem Brief an Christiaan Huygens\indexN{\textls{Huygens} (Hugenius, Ugenius, Hugens, Huguens), Christiaan 1629\textendash1695} vom 3./13.\;Oktober 1690 teilt er mit, eine bessere Methode für die Bestimmung der Tangenten durch \textit{foyers et filets} ausgearbeitet zu haben als diejenige Tschirnhausens,\indexN{\textls{Tschirnhaus} (H. v. Tch., H. v. Tsch.), Ehrenfried Walther v. 1651\textendash1708} wenngleich Fatio de Duillier\indexN{\textls{Fatio de Duillier} (Facius Dullerius), Nicolas 1664\textendash1753} bei der Veröffentlichung seiner eigenen Ergebnisse schneller gewesen sei als er sebst (\textit{LSB} III,\;4 N.\;283, S.\;623.17--19).\cite{01506}
Sodann erklärt er Huygens\indexN{\textls{Huygens} (Hugenius, Ugenius, Hugens, Huguens), Christiaan 1629\textendash1695} seinen methodischen Ansatz (ebd. S.\;623.21--624.5)\cite{01506} und stellt schließlich eine Betrachtung an, die eine bemerkenswerte Übereinstimmung mit der Ausführung auf \refpassage{LH_35_11_14_015_015r_bve-1}{LH_35_11_14_015_015r_bve-2} in der vorliegenden Aufzeichnung N.\;\ref{58943} aufweist: \textit{Cela posé je consideray que le stile qui tend les filets, peut estre conçu comme ayant autant de directions (egales en vistesse entre elles) qu'il y a de filets. Car comme il les tire, il en est tiré. Ainsi sa direction composée, qui doit estre dans la perpendiculaire à la courbe, passe par le centre de gravité d’autant de points qu’il y a de filets; qui sont les intersections d'un cercle (décrit du point de la courbe) avec les filets} (\textit{LSB} III,\;4 N.\;283, S.\;624.5--10).\cite{01506}
Es ist demgemäß unwahrscheinlich, dass Leibniz vor dem Herbst 1690 -- oder sogar vor seiner Rückkehr nach Hannover\indexO{Hannover} Mitte Juni 1690 -- Fatio de Duilliers\indexN{\textls{Fatio de Duillier} (Facius Dullerius), Nicolas 1664\textendash1753} \textit{Réflexions}\cite{02539} und \textit{Réponse}\cite{02540} gelesen hat.
Der Brief an Jacob Bernoulli\indexN{\textls{Bernoulli}, Jacob 1655\textendash1705}\cite{01260} und insbesondere derjenige an Christiaan Huygens\indexN{\textls{Huygens} (Hugenius, Ugenius, Hugens, Huguens), Christiaan 1629\textendash1695}\cite{01506} lassen vielmehr vermuten, dass die vorliegende Aufzeichnung N.\;\ref{58943} ungefähr zur gleichen Zeit entstanden sein dürfte.
Hieraus ergibt sich die für die Aufzeichnung N.\;\ref{58943} vorgeschlagene Datierung auf den Herbst 1690.
\pend%
\pstart%
Der Vollständigkeit halber ist festzuhalten, dass Leibniz bereits in seinem Brief an Rudolf Christian von Bodenhausen\indexN{\textls{Bodenhausen} (Bodenus, Bodenausen), Rudolf Christian, Freiherr v. 1698} vom 13.\;Januar 1690 die Methode zur Tangentenbestimmung dargelegt hatte, auf die er in den späteren Briefen an Jacob Bernoulli\indexN{\textls{Bernoulli}, Jacob 1655\textendash1705}\cite{01260} und vornehmlich an Christian Huygens\indexN{\textls{Huygens} (Hugenius, Ugenius, Hugens, Huguens), Christiaan 1629\textendash1695}\cite{01506} Bezug genommen hat (\textit{LSB} III,\;4 N.\;227, S.\;443.24--444.20)\cite{01385}.
Im Brief an Bodenhausen\indexN{\textls{Bodenhausen} (Bodenus, Bodenausen), Rudolf Christian, Freiherr v. 1698} erwähnt er freilich die unzulängliche Methode Tschirnhausens,\indexN{\textls{Tschirnhaus} (H. v. Tch., H. v. Tsch.), Ehrenfried Walther v. 1651\textendash1708} die er verbessern wollte, nicht jedoch -- anders als in den späteren Briefen an Bernoulli\indexN{\textls{Bernoulli}, Jacob 1655\textendash1705}\cite{01260} und Huygens\indexN{\textls{Huygens} (Hugenius, Ugenius, Hugens, Huguens), Christiaan 1629\textendash1695}\cite{01506} -- Fatio de Duilliers\indexN{\textls{Fatio de Duillier} (Facius Dullerius), Nicolas 1664\textendash1753} eigenständige Lösung.
Dies lässt sich als weiterer Beleg dafür betrachten, dass Leibniz Anfang 1690 die \textit{Réflexions}\cite{02539} und die \textit{Réponse}\cite{02540} noch nicht gelesen hatte.
Auch die Ausführung in dem während der Italienreise\indexO{Italien} entstandenen und im vorliegenden Band edierten Text N.\;\ref{58220} %
zeigt vielmehr, dass Leibniz wohl unabhängig vom Schweizer Mathematiker\indexN{\textls{Fatio de Duillier} (Facius Dullerius), Nicolas 1664\textendash1753} seine eigene Lösung entwickelt hatte.
Diese letztere machte er schließlich drei Jahre später in einem kurzen Aufsatz öffentlich: %
\textit{Deux problems construits par G.\:G.\;Leibniz en employant sa regle generale de la composition des mouvemens} (\textit{JS}, 14.\;September 1693, S.\;423\,f.; \textit{LMG} VI, S.\;233\,f.;\cite{00545} der Aufsatz wird voraussichtlich in einem späteren Band von \textit{LSB} VIII ediert werden).
\pend%
\pstart%
Es gilt schließlich zu bemerken, dass die in \textit{LSB} VI,\;4 N.\;382\cite{02546} edierten und auf \textit{Sommer bis Herbst 1688} (unter Vorbehalt) datierten Tschirnhaus-Auszüge\indexN{\textls{Tschirnhaus} (H. v. Tch., H. v. Tsch.), Ehrenfried Walther v. 1651\textendash1708} der oben vorgeschlagenen chronologischen Verortung der Aufzeichnung N.\;\ref{58943} nicht im Wege stehen.
In den Auszügen wird freilich eine Passage aus dem "Avertissement" zum Bd.\;V (April 1687) der \textit{Bibliothèque universelle et historique} (S.\;4\,r\textsuperscript{o}--5\,r\textsuperscript{o})\cite{02545} wörtlich wiedergegeben und mit der von Fatio de Duillier\indexN{\textls{Fatio de Duillier} (Facius Dullerius), Nicolas 1664\textendash1753} geäußerten Kritik an Tschirnhaus in Zusammenhang gesetzt (\textit{LSB} VI,\;4 N.\;382, S.\;2087.13--2088.11).\cite{02546}
Hieraus folgt jedoch nicht, dass Leibniz beim Verfassen dieser Auszüge Bd.\;V der \textit{Bibliothèque universelle et historique}\cite{02545} oder gar Fatio de Duilliers \textit{Réflexions}\cite{02539} vor Augen hatte.
Denn sämtliche in \textit{LSB} VI,\;4 N.\;382\cite{02546} edierten Auszüge entstammen entweder dem Aufsatz \textit{Eilfertiges Bedencken wieder die Objectiones, so im Mense Martio Schertz und Ernsthafter Gedancken über den Tractat Medicinae Mentis enthalten}, den Tschirnhaus in der Zeitschrift \textit{Freymüthige Lustige und Ernsthaffte jedoch Vernunfft- und Gesetz-Mässige Gedancken} (Juni 1688, S.\;746--792, bes. S.\;754\,f. u.\;762)\cite{02547} veröffentlicht hatte, oder der umfangreichen Erwiderung, die Christian Thomasius,\indexN{\textls{Thomasius}, Christian 1655\textendash1728} Herausgeber der Zeitschrift, im Anschluss an Tschirnhausens Aufsatz gedruckt hatte (ebd. S.\;793-854, bes. S.\;811--813, 814, 816\,f., 828, u.\;838\,f.).\cite{02548}
Ob die Auszüge \textit{LSB} VI,\;4 N.\;382\cite{02546} wirklich auf \textit{Sommer bis Herbst 1688} datiert werden sollen oder eher (wie es plausibel erscheint) etwas später entstanden, ist im Hinblick auf die Datierung der vorliegenden Aufzeichnung N.\;\ref{58943} unerheblich.
Denn sie setzen keineswegs voraus, dass Leibniz zum Zeitpunkt ihrer Entstehung direkte Kenntnis von Fatio de Duilliers\indexN{\textls{Fatio de Duillier} (Facius Dullerius), Nicolas 1664\textendash1753} Aufsätzen aus den Jahren 1687\cite{02539} und 1689\cite{02540} hatte, auf welche er in N.\;\ref{58943} anspielt.
\pend%
\pstart%
Ebenfalls unerheblich für die Datierung der vorliegenden Aufzeichnung N.\;\ref{58943} ist die Frage, wann das in \textit{LSB} VI,\;4 N.\;382,\cite{02546} S.\;2086.10--12 erwähnte Manuskript LH~XXXV\;12,\;1 Bl.\;42--43 verfasst wurde.
Dieses überliefert Leibnizens eigenhändige Auszüge der in der \textit{Bibliothèque universelle et historique} (Bd.\;III, September 1686, S.\;367--394)\cite{02545} veröffentlichten anonymen Besprechung\cite{02550} von Tschirnhausens Traktat \textit{Medicina mentis},\cite{02542} welcher bereits 1686, vor seinem laut Titelblatt offiziellen Erscheinungsdatum 1687, kursierte.
Diese Rezension kann für Leibniz aus offensichtlichen Gründen nicht die Quelle zu Fatio de Duilliers\indexN{\textls{Fatio de Duillier} (Facius Dullerius), Nicolas 1664\textendash1753} Kritik an Tschirnhaus\indexN{\textls{Tschirnhaus} (H. v. Tch., H. v. Tsch.), Ehrenfried Walther v. 1651\textendash1708} gewesen sein.
\pend%
\selectlanguage{latin}%
\frenchspacing%
\normalsize%
